% Candidate rules for change and cancellation.  Written by hand; see docs/one-order-coverage.md.

\section{Change and cancellation}

\begin{definition}[Cancelling an order (C-15)]
  \label{def:c15-cancel-order}
  \lean{Rules.OneOrder.Change.requestCancellation, Rules.OneOrder.Change.cancelDirect, Rules.OneOrder.Change.CancelTerms, Rules.OneOrder.Change.cancellationQuote}
  \uses{def:c06-item-status, def:payment, def:totals}
  An unpaid order can be cancelled directly.  A paid order is cancelled only by
  accepting a cancellation offer from the airline.  The model cancels an order
  at once when no payment has been received for it, cancelling and removing
  every service.  Otherwise it prices a cancellation offer from the order's
  price, as a change whose new price is zero.

  Source: \href{https://standards.atlassian.net/wiki/spaces/EASD/pages/574783569/Cancellation+of+Order+Items}{EASD guide, Cancellation of Order Items}.

  Capabilities: ORDCAN, ORDCA2.
\end{definition}

\begin{theorem}[Unpaid orders are cancelled at once]
  \label{thm:change-cancellation-unpaid-at-once}
  \lean{Rules.OneOrder.Change.unpaid_cancelled_at_once}
  \uses{def:c15-cancel-order}
  When no payment has been received for an order, a request to cancel it
  cancels and removes every service at once, and the order keeps its items.

  Source: \href{https://standards.atlassian.net/wiki/spaces/EASD/pages/574783569/Cancellation+of+Order+Items}{EASD guide, Cancellation of Order Items}.
\end{theorem}
\begin{proof}
  \uses{def:c15-cancel-order}
  Unfold the rule: with nothing paid it cancels the order directly.
\end{proof}

\begin{theorem}[Orders cancelled at once are closed]
  \label{thm:change-cancellation-cancelled-closed}
  \lean{Rules.OneOrder.Change.cancelled_at_once_closed}
  \uses{def:c15-cancel-order, def:c06-item-status, def:c07-order-status}
  Under every reading of the guide, an order cancelled at once has every item
  \texttt{CANCELLED} and is \texttt{CLOSED}.

  Source: \href{https://standards.atlassian.net/wiki/spaces/EASD/pages/574783569/Cancellation+of+Order+Items}{EASD guide, Cancellation of Order Items}.
\end{theorem}
\begin{proof}
  \uses{def:c15-cancel-order, def:c05-holds-value}
  Every service is cancelled and removed, and a removed service holds no value
  under any reading.
\end{proof}

\begin{theorem}[Penalty netted or charged separately (C-16)]
  \label{thm:c16-penalty-net}
  \lean{Rules.OneOrder.Change.net_same_either_way}
  \uses{def:c15-cancel-order, def:money, def:change-cancellation-due-amounts}
  A cancellation penalty is either deducted from the refund or charged
  separately, and the customer's net amount is the same either way.  For
  example, with 625.00 refundable and a penalty of 100.00, the customer is
  525.00 better off in both cases.  The model proves the two ways equal for
  changes and cancellations, under both readings of a penalty larger than the
  refundable total (C-85).

  Source: \href{https://standards.atlassian.net/wiki/spaces/EASD/pages/574783569/Cancellation+of+Order+Items}{EASD guide, Cancellation of Order Items}.

  Capabilities: ORDPEN.
\end{theorem}
\begin{proof}
  \uses{def:change-cancellation-due-amounts}
  Both ways, what the customer pays less what the airline returns is the price
  difference plus the penalty charged.
\end{proof}

\begin{definition}[Changing an order item (C-18)]
  \label{def:c18-change}
  \lean{Rules.OneOrder.Change.changeItem, Rules.OneOrder.Change.ItemChange, Rules.OneOrder.Change.originalAfterChange, Rules.OneOrder.Change.afterChange, Rules.OneOrder.Change.newItemOf}
  \uses{def:c01-structure, def:c06-item-status}
  Changing an order item cancels the original item's services, creates a new
  item, and charges or refunds the difference in price.  The model covers the
  change once it is paid: the services the change keeps move to the new item
  under the same references, their copies in the original item are marked
  \texttt{TRANSFERRED} with their delivery status unchanged, and the other
  services are cancelled and removed.  The money is priced by the price
  differential.

  Open point: the source shows the difference only in a worked example, and
  does not define its sign or how penalties enter it.

  Source: \href{https://standards.atlassian.net/wiki/spaces/EASD/pages/574619782/Modification+of+an+Order}{EASD guide, Modification of an Order}.

  Capabilities: ORDRSH.
\end{definition}

\begin{theorem}[Kept services keep their references]
  \label{thm:change-cancellation-kept-references}
  \lean{Rules.OneOrder.Change.change_keeps_services}
  \uses{def:c18-change}
  Every service a change keeps has a copy in the new item with the same
  reference, passenger, segment and delivery status.

  Source: \href{https://standards.atlassian.net/wiki/spaces/EASD/pages/574619782/Modification+of+an+Order}{EASD guide, Modification of an Order}.
\end{theorem}
\begin{proof}
  \uses{def:c18-change}
  The new item takes the kept services of the changed item and confirms them.
\end{proof}

\begin{theorem}[Originals of kept services are transferred]
  \label{thm:change-cancellation-originals-transferred}
  \lean{Rules.OneOrder.Change.change_marks_transferred}
  \uses{def:c18-change}
  Every service a change keeps stays in the changed item, marked
  \texttt{TRANSFERRED}, with its delivery status unchanged.

  Source: \href{https://standards.atlassian.net/wiki/spaces/EASD/pages/574619782/Modification+of+an+Order}{EASD guide, Modification of an Order}.
\end{theorem}
\begin{proof}
  \uses{def:c18-change}
  The changed item marks each service it keeps as transferred and changes
  nothing else about it.
\end{proof}

\begin{theorem}[Services not kept are cancelled]
  \label{thm:change-cancellation-others-cancelled}
  \lean{Rules.OneOrder.Change.change_cancels_others}
  \uses{def:c18-change}
  Every service of the changed item that the change does not keep is
  cancelled and removed.

  Source: \href{https://standards.atlassian.net/wiki/spaces/EASD/pages/574619782/Modification+of+an+Order}{EASD guide, Modification of an Order}.
\end{theorem}
\begin{proof}
  \uses{def:c18-change}
  The changed item marks each service it does not keep as cancelled and
  removed.
\end{proof}

\begin{theorem}[Changes keep item references distinct]
  \label{thm:change-cancellation-refs-distinct}
  \lean{Rules.OneOrder.Change.change_item_refs_unique}
  \uses{def:c18-change, thm:c03-unique-item-ids}
  If an order's item references are distinct before a change, they are
  distinct after it.  A change whose new item would reuse a reference is
  refused.

  Source: \href{https://standards.atlassian.net/wiki/spaces/EASD/pages/574619782/Modification+of+an+Order}{EASD guide, Modification of an Order}.
\end{theorem}
\begin{proof}
  \uses{def:c18-change}
  The changed item keeps its reference, and the new item's reference is not
  already in the order.
\end{proof}

\begin{theorem}[Transferred services keep the order balanced (C-19)]
  \label{thm:c19-transfer-balanced}
  \uses{def:c18-change, def:c06-item-status}
  \notready
  Services kept through a change move to the new item with the same service
  ID, and the originals are marked \texttt{TRANSFERRED}.  Delivered services
  move too, so that the order stays balanced.

  Open point: the source does not define what keeps an order balanced.

  Source: \href{https://standards.atlassian.net/wiki/spaces/EASD/pages/574849062/Order}{EASD guide, Order}.

  Capabilities: ORDRSH.
\end{theorem}

\begin{theorem}[Stored value stays in its order (C-20)]
  \label{thm:c20-value-store}
  \uses{def:c18-change, def:c05-holds-value}
  \notready
  Value stored as an order item can only be used in a later change to the same
  order.

  Open point: the source does not cover using part of the stored value, or
  when it expires.

  Source: \href{https://standards.atlassian.net/wiki/spaces/EASD/pages/574554245/Value+Store}{EASD guide, Value Store}.

  Capabilities: ORDREU.
\end{theorem}

% Second survey round: candidates written from the survey notes in docs/one-order-survey.md.

\begin{definition}[Price differential of a change (C-81)]
  \label{def:change-cancellation-price-differential}
  \lean{Rules.OneOrder.Change.priceDifferential, Rules.OneOrder.Change.ChangeQuote, Rules.OneOrder.Change.PriceDifferential, Rules.OneOrder.Change.DifferentialType, Rules.OneOrder.Change.statedDue, Rules.OneOrder.Change.typeOf}
  \uses{def:money, def:c18-change, def:c15-cancel-order, def:change-cancellation-due-amounts}
  Every change or cancellation offer carries a price differential: the old
  price, the new price (zero for a cancellation), the penalties, the
  forfeited amounts, and two non-negative amounts, one due to the airline
  and one due by the airline. A type code names the outcome: EvenExchange,
  AddCol, Refund, Residual (a value store), Reusable, AddColAndRefund or
  AddColAndResidual.

  Open point: The sign of the differential total is inconsistent. The
  Cancellation of Order Items page shows a refund with a positive total of
  625.00, while the worked examples show refunds with negative totals, such
  as minus 1,000.00 in \texttt{EXM\_ORD\_030A}. The business must fix one
  sign rule, or model only the two non-negative due amounts. The model states
  only the two due amounts until the business decides.

  Sources: \href{https://standards.atlassian.net/wiki/spaces/EASD/pages/574783569}{EASD guide, Cancellation of Order Items};
  \href{https://standards.atlassian.net/wiki/spaces/EASD/pages/842268673}{EASD guide, Customer requests a change to an Order containing a paid OrderItem};
  \href{https://standards.atlassian.net/wiki/spaces/EASD/pages/817594369}{EASD guide, Customer requests a cancellation to paid Order or OrderItems};
  \href{https://retailing.iata.org/armi/docs/ORDRSH/}{IATA capability ORDRSH};
  \href{https://standards.atlassian.net/wiki/spaces/EASD/pages/574554213/Price+Differentials}{EASD guide, Price Differentials}.

  Capabilities: ORDRSH, ORDCAN, ORDCA2, ORDPEN, ORDFOR, ORDREU.
\end{definition}

\begin{definition}[Amounts due to and by the airline]
  \label{def:change-cancellation-due-amounts}
  \lean{Rules.OneOrder.Change.dueAmounts, Rules.OneOrder.Change.DueAmounts, Rules.OneOrder.Change.ChangeQuote.difference, Rules.OneOrder.Change.ChangeQuote.refundable, Rules.OneOrder.Change.ExcessPenalty.charged, Rules.OneOrder.Change.PenaltyHandling, Rules.OneOrder.Change.Settlement}
  \uses{def:money, def:change-cancellation-penalty-exceeds-refund}
  The price difference of a change is the new price, less the old price, plus
  the forfeited amount.  With a netted penalty, the penalty charged is added
  to the difference, and the result is due to the airline when positive and
  by the airline when negative.  With a separate penalty, the penalty charged
  is due to the airline on its own, and the difference is settled on its own.
  The worked examples fit this formula; the guide does not state it.

  Source: \href{https://standards.atlassian.net/wiki/spaces/EASD/pages/574554213/Price+Differentials}{EASD guide, Price Differentials}.
\end{definition}

\begin{theorem}[Mixed currencies are rejected]
  \label{thm:change-cancellation-mixed-currency}
  \lean{Rules.OneOrder.Change.differential_rejects_mixed_currency}
  \uses{def:change-cancellation-price-differential}
  A change whose prices, penalty and forfeited amount are not all in one
  currency is rejected unpriced.

  Source: \href{https://standards.atlassian.net/wiki/spaces/EASD/pages/574619782/Modification+of+an+Order}{EASD guide, Modification of an Order}.
\end{theorem}
\begin{proof}
  \uses{def:change-cancellation-price-differential}
  The rule checks the currencies first.
\end{proof}

\begin{theorem}[Type code names a collection]
  \label{thm:change-cancellation-type-collects}
  \lean{Rules.OneOrder.Change.differential_collects_iff}
  \uses{def:change-cancellation-price-differential}
  A differential's type code is AddCol, AddColAndRefund or AddColAndResidual
  exactly when an amount is due to the airline.

  Source: \href{https://standards.atlassian.net/wiki/spaces/EASD/pages/574554213/Price+Differentials}{EASD guide, Price Differentials}.
\end{theorem}
\begin{proof}
  \uses{def:change-cancellation-price-differential}
  Check each branch of the type code against the due amounts.
\end{proof}

\begin{theorem}[Type code names a refund or a value store]
  \label{thm:change-cancellation-type-pays}
  \lean{Rules.OneOrder.Change.differential_pays_iff}
  \uses{def:change-cancellation-price-differential}
  A differential's type code is Refund, Residual, AddColAndRefund or
  AddColAndResidual exactly when an amount is due by the airline.

  Source: \href{https://standards.atlassian.net/wiki/spaces/EASD/pages/574554213/Price+Differentials}{EASD guide, Price Differentials}.
\end{theorem}
\begin{proof}
  \uses{def:change-cancellation-price-differential}
  Check each branch of the type code against the due amounts.
\end{proof}

\begin{theorem}[Value kept for later use is typed Reusable]
  \label{thm:change-cancellation-respend-reusable}
  \lean{Rules.OneOrder.Change.respend_reusable}
  \uses{def:change-cancellation-price-differential, def:change-cancellation-reusable-amount}
  When the seller asks to keep the value of cancelled items for later use, the
  differential is typed Reusable and both due amounts are zero.

  Source: \href{https://standards.atlassian.net/wiki/spaces/EASD/pages/817594921}{EASD guide, Customer requests a cancellation to paid Order or OrderItems resulting in a Reusable amount}.
\end{theorem}
\begin{proof}
  \uses{def:change-cancellation-price-differential}
  Unfold the rule: a request to keep the value states both amounts as zero.
\end{proof}

\begin{theorem}[Netted penalty gives one amount]
  \label{thm:change-cancellation-netted-one-amount}
  \lean{Rules.OneOrder.Change.netted_single_amount}
  \uses{def:change-cancellation-due-amounts}
  With a netted penalty, the due amounts net to the price difference plus the
  penalty charged, and at most one of them is above zero.

  Source: \href{https://standards.atlassian.net/wiki/spaces/EASD/pages/574783569/Cancellation+of+Order+Items}{EASD guide, Cancellation of Order Items}.
\end{theorem}
\begin{proof}
  \uses{def:change-cancellation-due-amounts}
  A netted result is split by its sign.
\end{proof}

\begin{theorem}[Separate penalty is added to the amount due]
  \label{thm:change-cancellation-separate-added}
  \lean{Rules.OneOrder.Change.separate_penalty_added}
  \uses{def:change-cancellation-due-amounts}
  With a separate penalty, the amount due to the airline is the penalty
  charged plus any price difference the customer pays, and the amount due by
  the airline is any price difference the airline returns.

  Source: \href{https://standards.atlassian.net/wiki/spaces/EASD/pages/574783569/Cancellation+of+Order+Items}{EASD guide, Cancellation of Order Items}.
\end{theorem}
\begin{proof}
  \uses{def:change-cancellation-due-amounts}
  Unfold the rule.
\end{proof}

\begin{theorem}[Due amounts net to the difference (C-82)]
  \label{thm:change-cancellation-due-net-identity}
  \lean{Rules.OneOrder.Change.differential_nets}
  \uses{def:change-cancellation-price-differential, thm:c16-penalty-net, def:c18-change, def:money}
  For every change or cancellation offer, the amount due to the airline
  minus the amount due by the airline equals the new price minus the old
  price, plus penalties, plus forfeited amounts. Netting of penalties or
  taxes changes how this net is split between the two due amounts and never
  changes the net. In \texttt{EXM\_ORD\_003C}, 900 minus 1,000 plus a 40
  penalty plus 120 forfeited gives 60 due to the airline.

  Case to check: The sample change record on the Modification of an Order
  page shows a Refund with old and new prices of zero and 100 due by the
  airline, in currency XSD against prices in USD; no penalty or forfeit
  explains it. A Reusable cancellation (\texttt{EXM\_ORD\_030E}) shows every
  amount as zero, so the identity holds only if the value kept on the
  cancelled item counts as the new price.

  Counterexample: a paid order at 1,000.00 cancelled with a request to keep
  the value for later use is typed Reusable with both amounts zero, where the
  formula gives 1,000.00 due by the airline (Lean example
  \texttt{respendCancellation}, after \texttt{EXM\_ORD\_030E}).  Under the
  reading that waives a penalty above the refundable value, a cancellation of
  1,000.00 with 20.00 forfeited and a 1,000.00 penalty has nothing due, where
  the formula gives 20.00 due to the airline (\texttt{penaltyAboveRefund}).

  Open point: The equation holds in all eleven worked examples that settle
  money, but the guide never states it, and C-18 leaves the sign and the
  penalty term open. The business must adopt this formula, or another one,
  as the definition of the amount owed on a change.

  Sources: \href{https://standards.atlassian.net/wiki/spaces/EASD/pages/842268745}{EASD guide, Customer requests a change of a paid order prior to travel resulting in additional collection due to penalty and forfeited amount};
  \href{https://standards.atlassian.net/wiki/spaces/EASD/pages/842268768}{EASD guide, Customer requests a change of a paid order prior to travel resulting in a refund};
  \href{https://standards.atlassian.net/wiki/spaces/EASD/pages/842301447}{EASD guide, Customer requests a change of a paid order prior to travel resulting in additional collection and a refund};
  \href{https://standards.atlassian.net/wiki/spaces/EASD/pages/817594543}{EASD guide, Customer requests a cancellation to paid Order or OrderItems resulting in Partial Refund with a Forfeited amount incurred};
  \href{https://standards.atlassian.net/wiki/spaces/EASD/pages/817594795}{EASD guide, Customer requests a cancellation to paid Order or OrderItems resulting in Partial Refund with a Forfeited amount incurred and requiring payment of a penalty};
  \href{https://standards.atlassian.net/wiki/spaces/EASD/pages/574619782}{EASD guide, Modification of an Order}.

  Capabilities: ORDRSH, ORDCAN, ORDPEN, ORDFOR.
\end{theorem}
\begin{proof}
  \uses{def:change-cancellation-price-differential}
  Not proven.  The counterexample above shows the statement is false for the
  rule as written.
\end{proof}

\begin{theorem}[Collecting the whole penalty keeps the formula]
  \label{thm:change-cancellation-net-of-collect}
  \lean{Rules.OneOrder.Change.due_amounts_net_of_collect}
  \uses{def:change-cancellation-due-amounts, def:change-cancellation-penalty-exceeds-refund}
  Under the reading that collects the whole penalty, the due amounts net to
  the new price, less the old price, plus the penalty, plus the forfeited
  amount, whether the penalty is netted or charged separately.

  Source: \href{https://standards.atlassian.net/wiki/spaces/EASD/pages/574783569/Cancellation+of+Order+Items}{EASD guide, Cancellation of Order Items}.
\end{theorem}
\begin{proof}
  \uses{def:change-cancellation-due-amounts}
  The penalty charged is the whole penalty, and either way the net is the
  price difference plus the penalty charged.
\end{proof}

\begin{theorem}[Order balance after servicing (C-83)]
  \label{thm:change-cancellation-order-balance}
  \lean{Rules.OneOrder.Change.change_keeps_balance}
  \uses{def:totals, def:payment, def:refund, thm:c19-transfer-balanced, def:c06-item-status, thm:change-cancellation-penalties-persist, thm:change-cancellation-due-net-identity, def:change-cancellation-ledger, def:c18-change}
  At every point in an order's life, payments minus refunds equal the prices
  of its live order items plus all recorded penalties plus all recorded
  forfeited amounts. Value store and reusable items count as live at their
  price. In \texttt{EXM\_ORD\_030B}, 1,000 paid minus 930 refunded equals a
  50 penalty plus 20 forfeited.

  Case to check: In \texttt{EXM\_ORD\_003C} after acceptance and before
  payment, the old item is still \texttt{ACTIVE} at 1,000 and the new item
  is \texttt{ACTIVE} at 900 with 1,000 paid, so the equation fails until the
  old item stops counting. In the Modification of an Order illustration,
  payments total 3,100 while the new item shows a price of 3,000 against a
  quoted new price of 3,100.

  Counterexample: a paid trip at 1,000.00 changed to a trip at 900.00 that
  keeps the outbound flight, with a 40.00 penalty and 120.00 forfeited, leaves
  1,060.00 paid against 1,900.00 of live items plus 160.00 recorded, because
  the original item still holds value.  The same change keeping nothing
  balances (Lean example \texttt{changeBalance}).

  Open point: The ORDPEN and ORDFOR guidance requires that sellers can
  reconcile all order values against payments and refunds but gives no
  equation, and C-19 leaves balance undefined. The business must choose the
  equation and say how an accepted but unpaid change enters it.

  Sources: \href{https://retailing.iata.org/armi/docs/ORDPEN/}{IATA capability ORDPEN};
  \href{https://retailing.iata.org/armi/docs/ORDFOR/}{IATA capability ORDFOR};
  \href{https://standards.atlassian.net/wiki/spaces/EASD/pages/842268745}{EASD guide, Customer requests a change of a paid order prior to travel resulting in additional collection due to penalty and forfeited amount};
  \href{https://standards.atlassian.net/wiki/spaces/EASD/pages/817594543}{EASD guide, Customer requests a cancellation to paid Order or OrderItems resulting in Partial Refund with a Forfeited amount incurred};
  \href{https://standards.atlassian.net/wiki/spaces/EASD/pages/574619782}{EASD guide, Modification of an Order};
  \href{https://retailing.iata.org/armi/docs/PAYREF/}{IATA capability PAYREF};
  \href{https://standards.atlassian.net/wiki/spaces/EASD/pages/842268768/Customer+requests+a+change+of+a+paid+order+prior+to+travel+resulting+in+a+refund}{EASD guide, Customer requests a change of a paid order prior to travel resulting in a refund};
  \href{https://standards.atlassian.net/wiki/spaces/EASD/pages/842301447/Customer+requests+a+change+of+a+paid+order+prior+to+travel+resulting+in+additional+collection+and+a+refund}{EASD guide, Customer requests a change of a paid order prior to travel resulting in additional collection and a refund};
  \href{https://standards.atlassian.net/wiki/spaces/EASD/pages/817594417/Customer+requests+a+cancellation+to+paid+Order+or+OrderItems+resulting+in+full+refund}{EASD guide, Customer requests a cancellation to paid Order or OrderItems resulting in full refund}.

  Capabilities: ORDPEN, ORDFOR, ORDRSH, PAYREF, PAYSUM.
\end{theorem}
\begin{proof}
  \uses{def:change-cancellation-ledger, def:c18-change}
  Not proven.  The counterexample above shows the statement is false for the
  rule as written.
\end{proof}

\begin{definition}[Money recorded on an order]
  \label{def:change-cancellation-ledger}
  \lean{Rules.OneOrder.Change.settle, Rules.OneOrder.Change.Ledger, Rules.OneOrder.Change.liveValue, Rules.OneOrder.Change.Ledger.Balanced}
  \uses{def:money, def:change-cancellation-price-differential, def:c05-holds-value}
  An order records everything the customer paid, everything the airline
  returned, and every penalty and forfeited amount, at order level.  Settling
  a change or cancellation offer adds its amount due to the airline to the
  payments, its amount due by the airline to the refunds, and its penalty and
  forfeited amount to those recorded.  The live items are the order's active
  items.

  Source: \href{https://retailing.iata.org/armi/docs/ORDPEN/}{IATA capability ORDPEN}.
\end{definition}

\begin{theorem}[Settling adds to payments and refunds]
  \label{thm:change-cancellation-settle-payments}
  \lean{Rules.OneOrder.Change.settle_adds_payments}
  \uses{def:change-cancellation-ledger}
  Settling an offer adds the amount due to the airline to the payments and the
  amount due by the airline to the refunds.

  Source: \href{https://retailing.iata.org/armi/docs/ORDPEN/}{IATA capability ORDPEN}.
\end{theorem}
\begin{proof}
  \uses{def:change-cancellation-ledger}
  Unfold the rule.
\end{proof}

\begin{theorem}[Settling records penalties and forfeits]
  \label{thm:change-cancellation-settle-records}
  \lean{Rules.OneOrder.Change.settle_records_penalty}
  \uses{def:change-cancellation-ledger}
  Settling an offer adds its penalty and forfeited amount to those recorded on
  the order, and lowers no amount recorded before.

  Source: \href{https://retailing.iata.org/armi/docs/ORDPEN/}{IATA capability ORDPEN}.
\end{theorem}
\begin{proof}
  \uses{def:change-cancellation-ledger}
  Settling only adds.
\end{proof}

\begin{definition}[Change accepted before payment (C-84)]
  \label{def:change-cancellation-deferred-payment-change}
  \uses{def:c11-payment-time-limit, def:c08-delivery-status, def:c05-holds-value, def:c18-change}
  \notready
  A seller may accept a change offer and pay later. Until payment, the new
  item's services are \texttt{CONFIRMED} with a payment time limit, and the
  old item stays \texttt{ACTIVE} with its replaced services
  \texttt{CANCELLED} and in \texttt{REMOVED HOLD VALUE}. Payment moves the
  new services to \texttt{READY TO PROCEED} and the old item to
  \texttt{CANCELLED}.

  Case to check: In \texttt{EXM\_ORD\_003C} the change is accepted late on
  18 May 2023 with payment due by 16:00 UTC on 19 May. If the 60.00 is never
  paid, neither item gives the customer a flight.

  Open point: The source does not say what happens if the payment time limit
  passes. The business must decide whether the old services return to
  \texttt{READY TO PROCEED}, or the customer is left with an unpaid new item
  and cancelled old services.

  Sources: \href{https://standards.atlassian.net/wiki/spaces/EASD/pages/842268673}{EASD guide, Customer requests a change to an Order containing a paid OrderItem};
  \href{https://standards.atlassian.net/wiki/spaces/EASD/pages/842268745}{EASD guide, Customer requests a change of a paid order prior to travel resulting in additional collection due to penalty and forfeited amount};
  \href{https://standards.atlassian.net/wiki/spaces/EASD/pages/842301447}{EASD guide, Customer requests a change of a paid order prior to travel resulting in additional collection and a refund};
  \href{https://standards.atlassian.net/wiki/spaces/EASD/pages/842301447/Customer+requests+a+change+of+a+paid+order+prior+to+travel+resulting+in+additional+collection+and+a+refund}{EASD guide, Customer requests a change of a paid order prior to travel resulting in additional collection and a refund};
  \href{https://standards.atlassian.net/wiki/spaces/EASD/pages/842268768/Customer+requests+a+change+of+a+paid+order+prior+to+travel+resulting+in+a+refund}{EASD guide, Customer requests a change of a paid order prior to travel resulting in a refund}.

  Capabilities: ORDRSH, ORDPEN, ORDFOR, PAYREF, PAYORD.
\end{definition}

\begin{definition}[Penalty larger than the refund (C-85)]
  \label{def:change-cancellation-penalty-exceeds-refund}
  \lean{Rules.OneOrder.Change.ExcessPenalty}
  \uses{def:money, def:c15-cancel-order}
  With a netted penalty, the amount due by the airline is the refundable
  value of the cancelled items minus the penalty. The source has no case
  where the penalty is larger than the refundable value.

  Case to check: Take \texttt{EXM\_ORD\_030B} with a cancellation fee of
  1,000 in place of 50. The refundable value is 980, so the netted result is
  20 due to the airline on an offer typed Refund.

  Open point: The business must decide whether a penalty is capped at the
  refundable value, so that cancelling never leaves the customer owing
  money, or whether the excess is collected. The model takes the answer as a
  parameter with two readings: collect the excess, or waive it.

  Sources: \href{https://standards.atlassian.net/wiki/spaces/EASD/pages/574783569}{EASD guide, Cancellation of Order Items};
  \href{https://standards.atlassian.net/wiki/spaces/EASD/pages/817594543}{EASD guide, Customer requests a cancellation to paid Order or OrderItems resulting in Partial Refund with a Forfeited amount incurred}.

  Capabilities: ORDCAN, ORDCA2, ORDPEN.
\end{definition}

\begin{theorem}[Cancelling leaves the customer owing nothing]
  \label{thm:change-cancellation-cancel-not-owing}
  \lean{Rules.OneOrder.Change.cancel_not_owing}
  \uses{def:change-cancellation-due-amounts, def:change-cancellation-penalty-exceeds-refund}
  On a cancellation whose forfeited amount is part of the old price, the
  amount due to the airline is no larger than the amount due by the airline,
  under both readings of a penalty above the refundable value.

  Counterexample: a paid order at 1,000.00 with 20.00 forfeited, cancelled
  with a netted penalty of 1,000.00, leaves 20.00 due to the airline under the
  reading that collects the whole penalty (Lean example
  \texttt{penaltyAboveRefund}).

  Open point: under the other reading the guarantee holds, but the due
  amounts no longer net to the formula of C-82.  The business must decide
  whether a penalty is capped at the refundable value or the excess is
  collected.

  Source: \href{https://standards.atlassian.net/wiki/spaces/EASD/pages/574783569/Cancellation+of+Order+Items}{EASD guide, Cancellation of Order Items}.
\end{theorem}
\begin{proof}
  \uses{def:change-cancellation-due-amounts}
  Not proven.  The counterexample above shows the statement is false for the
  rule as written.
\end{proof}

\begin{theorem}[Waiving the excess penalty leaves nothing owing]
  \label{thm:change-cancellation-not-owing-of-waive}
  \lean{Rules.OneOrder.Change.cancel_not_owing_of_waive}
  \uses{def:change-cancellation-due-amounts, def:change-cancellation-penalty-exceeds-refund}
  Under the reading that waives a penalty above the refundable value, a
  cancellation never leaves the customer owing the airline money.

  Source: \href{https://standards.atlassian.net/wiki/spaces/EASD/pages/574783569/Cancellation+of+Order+Items}{EASD guide, Cancellation of Order Items}.
\end{theorem}
\begin{proof}
  \uses{def:change-cancellation-due-amounts}
  The penalty charged is at most the refundable value, and a cancellation
  returns the refundable value.
\end{proof}

\begin{theorem}[Full cancellation keeps no residual (C-86)]
  \label{thm:change-cancellation-full-cancel-no-residual}
  \uses{def:c15-cancel-order, thm:c20-value-store, def:c07-order-status}
  \notready
  Cancelling a whole order cancels every service and settles any refundable
  amount. The customer cannot keep residual value in the order for a later
  reshop.

  Case to check: If a full cancellation must close the order, 030C conflicts
  with C-07, because its value store service still holds value and keeps the
  order \texttt{OPENED}.

  Open point: The ORDCA2 guidance states this, but worked example
  \texttt{EXM\_ORD\_030C} cancels a whole order through the full
  cancellation request and keeps the full 1,000 as a Value Store order item,
  leaving the order \texttt{OPENED}. The business must decide whether a full
  cancellation may end in a value store or a reusable amount.

  Sources: \href{https://retailing.iata.org/armi/docs/ORDCA2/}{IATA capability ORDCA2};
  \href{https://standards.atlassian.net/wiki/spaces/EASD/pages/817594669}{EASD guide, Customer requests a cancellation to paid Order or OrderItems resulting in Value Store}.

  Capabilities: ORDCA2, ORDREU.
\end{theorem}

\begin{definition}[Paid or partly paid orders need an offer (C-87)]
  \label{def:change-cancellation-partly-paid-cancel}
  \lean{Rules.OneOrder.Change.paid_needs_offer}
  \uses{def:c15-cancel-order, thm:refunds-le-payments, def:payment}
  An order with any payment, full or partial, is cancelled through a
  cancellation offer. Only an order with no payment at all may be cancelled
  directly.

  Case to check: Take an order priced 1,000 with 300 paid. The worked
  examples compute refundable value from the old price, which would show
  1,000 due by the airline against 300 received and break the rule that
  refunds never exceed payments.

  Open point: The EASD page speaks only of paid and unpaid orders; the
  partial case comes from ORDCA2. The business must decide which route
  applies when some items are paid and others are not, and whether the
  refundable amount is capped at the amount received.

  Sources: \href{https://retailing.iata.org/armi/docs/ORDCA2/}{IATA capability ORDCA2};
  \href{https://standards.atlassian.net/wiki/spaces/EASD/pages/574783569}{EASD guide, Cancellation of Order Items}.

  Capabilities: ORDCA2, ORDCAN.
\end{definition}

\begin{theorem}[Cancellation offers stay within payments]
  \label{thm:change-cancellation-offer-within-payments}
  \lean{Rules.OneOrder.Change.offer_within_payments}
  \uses{def:c15-cancel-order, def:change-cancellation-partly-paid-cancel, def:change-cancellation-price-differential, def:totals}
  A cancellation offer never makes the amount due by the airline larger than
  the payments received for the order.

  Counterexample: an order at 1,000.00 with 300.00 received is cancelled
  through an offer priced from the order's price, with 1,000.00 due by the
  airline (Lean example \texttt{partlyPaidCancellation}).

  Open point: the business must decide whether the refundable amount is
  capped at the payments received, and which route applies when some items
  are paid and others are not.

  Source: \href{https://retailing.iata.org/armi/docs/ORDCA2/}{IATA capability ORDCA2}.
\end{theorem}
\begin{proof}
  \uses{def:c15-cancel-order, def:change-cancellation-price-differential}
  Not proven.  The counterexample above shows the statement is false for the
  rule as written.
\end{proof}

\begin{definition}[Removing one passenger from an item (C-88)]
  \label{def:change-cancellation-remove-passenger}
  \uses{thm:c04-item-unsplit, thm:c19-transfer-balanced, def:c18-change, def:passenger}
  \notready
  When one order item covers several passengers, removing a passenger
  deletes the whole item and lists the other passengers' services to retain.
  The retained services move to a new item, which is repriced.

  Case to check: Two passengers share one item priced 1,000 under a
  companion fare. Removing one reprices the other at a 900 single fare, so
  the leaving passenger is refunded 100 where an even split gives 500.

  Open point: The source does not say whether the retained services keep
  their share of the original price or are repriced at current fares, or how
  a penalty is assigned when one passenger leaves. The same page also
  classes cancelling services inside an item as a change of order, yet uses
  that method for passenger removal under cancellation.

  Sources: \href{https://standards.atlassian.net/wiki/spaces/EASD/pages/574783569}{EASD guide, Cancellation of Order Items};
  \href{https://retailing.iata.org/armi/docs/ORDCAN/}{IATA capability ORDCAN}.

  Capabilities: ORDCAN, ORDRSH.
\end{definition}

\begin{definition}[Reusable amount on cancellation (C-89)]
  \label{def:change-cancellation-reusable-amount}
  \uses{def:c05-holds-value, def:c06-item-status, def:c07-order-status, thm:c20-value-store}
  \notready
  A cancellation that asks to respend (RespendInd true) triggers no refund.
  The item stays \texttt{ACTIVE} at its full price, its services become
  \texttt{CANCELLED} with delivery status \texttt{REMOVED HOLD VALUE}, the
  differential is typed Reusable with zero amounts, and the order stays
  \texttt{OPENED}. The exact amount is computed later, at a refund request
  or a reshop.

  Case to check: In \texttt{EXM\_ORD\_030E} the customer paid 1,000 and
  holds an \texttt{ACTIVE} item at 1,000. If the airline recomputes the
  value lower at reuse, nothing in the source protects the 1,000.

  Open point: The source sets no floor, ceiling or time limit on the amount
  computed later. The business must decide whether the customer is owed at
  least what a refund would have paid on the day of cancellation, and when
  the held value lapses.

  Sources: \href{https://retailing.iata.org/armi/docs/ORDREU/}{IATA capability ORDREU};
  \href{https://standards.atlassian.net/wiki/spaces/EASD/pages/817594921}{EASD guide, Customer requests a cancellation to paid Order or OrderItems resulting in a Reusable amount};
  \href{https://standards.atlassian.net/wiki/spaces/EASD/pages/574783569}{EASD guide, Cancellation of Order Items}.

  Capabilities: ORDREU, ORDCAN.
\end{definition}

\begin{definition}[Settling money the airline owes (C-90)]
  \label{def:change-cancellation-owed-settlement}
  \uses{def:refund, def:c17-refund-methods, thm:c20-value-store, def:change-cancellation-reusable-amount, def:change-cancellation-price-differential}
  \notready
  Money the airline owes after a change or cancellation is settled as a
  refund (by default to the original form of payment), a Value Store order
  item, a reusable amount held on the cancelled item, or value on an open
  accountable document. Non-refundable parts, such as a non-refundable tax,
  are recorded as forfeited.

  Case to check: In \texttt{EXM\_ORD\_003D} and 003E the refund offer and
  the value store offer are both worth 150. Nothing in the source stops a
  150 refund offer alongside a 120 value store offer.

  Open point: The sources differ on who picks the form. The change page has
  the airline choose between refund and value store, the cancellation page
  lets the seller request respend, and \texttt{EXM\_ORD\_003D} and 003E show
  the airline offering both while the seller accepts one. The business must
  decide who chooses, and whether a value store must equal the refund it
  replaces.

  Sources: \href{https://standards.atlassian.net/wiki/spaces/EASD/pages/574554168}{EASD guide, Seller Initiated Change to an Order};
  \href{https://standards.atlassian.net/wiki/spaces/EASD/pages/574783569}{EASD guide, Cancellation of Order Items};
  \href{https://standards.atlassian.net/wiki/spaces/EASD/pages/574554245}{EASD guide, Value Store};
  \href{https://retailing.iata.org/armi/docs/ORDCAN/}{IATA capability ORDCAN};
  \href{https://standards.atlassian.net/wiki/spaces/EASD/pages/842268768}{EASD guide, Customer requests a change of a paid order prior to travel resulting in a refund};
  \href{https://standards.atlassian.net/wiki/spaces/EASD/pages/842268791}{EASD guide, Customer requests a change of a paid order prior to travel resulting in a value store}.

  Capabilities: ORDCAN, ORDRSH, ORDREU, ORDFOR.
\end{definition}

\begin{definition}[Value store time windows (C-91)]
  \label{def:change-cancellation-value-store-windows}
  \uses{thm:c20-value-store, def:instant, def:c05-holds-value, def:c07-order-status}
  \notready
  A Value Store order item carries its own cancel and change restrictions
  with start and end dates. In \texttt{EXM\_ORD\_030C} the 1,000 value store
  may be cancelled at no fee for one year and used in a change for two
  years, and cancelling is not allowed after the first year.

  Case to check: A customer who holds the 030C value store for 25 months can
  neither cancel nor use it, yet the order stays \texttt{OPENED} with a
  service that still holds value under C-05.

  Open point: The source does not say what the customer receives on
  cancelling a value store, or what happens to unused value when the change
  window closes: refund, forfeit or a voucher. C-20 leaves expiry open;
  these windows supply the dates but not the outcome.

  Sources: \href{https://standards.atlassian.net/wiki/spaces/EASD/pages/817594669}{EASD guide, Customer requests a cancellation to paid Order or OrderItems resulting in Value Store};
  \href{https://standards.atlassian.net/wiki/spaces/EASD/pages/842268791}{EASD guide, Customer requests a change of a paid order prior to travel resulting in a value store};
  \href{https://standards.atlassian.net/wiki/spaces/EASD/pages/574554245}{EASD guide, Value Store}.

  Capabilities: ORDCA2, ORDCAN.
\end{definition}

\begin{theorem}[Stored value is used once (C-92)]
  \label{thm:change-cancellation-stored-value-used-once}
  \uses{thm:c20-value-store, def:c05-holds-value, def:change-cancellation-owed-settlement}
  \notready
  The total drawn from any stored value, whether a Value Store item, a
  reusable item or an open accountable document, never exceeds its value
  across all orders. Value moved to a new order leaves the order it came
  from.

  Case to check: A seller reallocates 500 from an open ticket to a new order
  through OrderCreate, then quotes a change on the original order naming the
  same ticket. Nothing in the source marks the 500 as spent.

  Open point: The Value Store page lets value on an open document be
  reallocated to a new order or within the changed order, recalculated on
  reuse, but states no rule against using it twice. The business must decide
  how stored value is tracked and when it is used up, including the partial
  use that C-20 leaves open.

  Source: \href{https://standards.atlassian.net/wiki/spaces/EASD/pages/574554245}{EASD guide, Value Store}.

  Capabilities: ORDREU.
\end{theorem}

\begin{theorem}[Transferred originals hold no value (C-93)]
  \label{thm:change-cancellation-transferred-original-no-value}
  \lean{Rules.OneOrder.Change.change_original_cancelled}
  \uses{def:c05-holds-value, def:c06-item-status, thm:c19-transfer-balanced, def:c08-delivery-status, def:c18-change}
  A service transferred to a new item holds value only in the new item. The
  original copy, with status \texttt{TRANSFERRED}, holds no value, so its
  old item can be \texttt{CANCELLED}.

  Case to check: In \texttt{EXM\_ORD\_003B} after the change, ORDITM01 is
  \texttt{CANCELLED} while six of its services show \texttt{TRANSFERRED}
  with \texttt{READY TO PROCEED}. The same six service IDs appear in
  \texttt{ORDITM02\_NEW} as \texttt{READY TO PROCEED}, so their value is
  counted twice.

  Counterexample: a paid return trip changed to keep the outbound flight
  leaves the original item with the outbound service \texttt{TRANSFERRED} and
  \texttt{READY TO PROCEED}.  Under both readings the original item is
  \texttt{ACTIVE}, and the outbound service holds value in both items (Lean
  example \texttt{changeKeepingOutbound}).

  Open point: In the worked examples the original copy keeps delivery status
  \texttt{READY TO PROCEED} while its item is \texttt{CANCELLED}. C-05
  judges value by delivery status alone, under which the cancelled item
  still holds value and meets the \texttt{ACTIVE} condition of C-06. The
  business must decide whether the \texttt{TRANSFERRED} status overrides the
  delivery status.

  Sources: \href{https://standards.atlassian.net/wiki/spaces/EASD/pages/841187358}{EASD guide, Customer requests a change of a paid order prior to travel resulting in additional collection};
  \href{https://retailing.iata.org/armi/docs/ORDRSH/}{IATA capability ORDRSH};
  \href{https://standards.atlassian.net/wiki/spaces/EASD/pages/574849062/Order}{EASD guide, Order};
  \href{https://standards.atlassian.net/wiki/spaces/OA/pages/844202130}{Order Accounting, Customer requests a change of a paid order prior to travel resulting in additional collection and a refund};
  \href{https://standards.atlassian.net/wiki/spaces/EASD/pages/842268712}{EASD guide, Customer requests a change of a paid order prior to travel resulting in additional collection and a value store}.

  Capabilities: ORDRSH.
\end{theorem}
\begin{proof}
  \uses{def:c18-change, def:c06-item-status}
  Not proven.  The counterexample above shows the statement is false for the
  rule as written.
\end{proof}

\begin{theorem}[Refund waits for separate penalty (C-94)]
  \label{thm:change-cancellation-refund-after-penalty}
  \uses{thm:c16-penalty-net, def:refund, def:payment, thm:c20-value-store}
  When a penalty is charged separately, the airline releases the refund, or
  makes the value store usable, only after the penalty is paid. A netted
  penalty leaves nothing to wait for.

  Case to check: C-16 treats netted and separate penalties as equivalent. In
  \texttt{EXM\_ORD\_003G} the separate penalty means the customer pays 60 in
  cash and receives a 160 value store; netted, the customer would pay
  nothing and receive 100 stored. The net is equal, but 60 of cash becomes
  restricted stored value.

  Sources: \href{https://standards.atlassian.net/wiki/spaces/EASD/pages/817594369}{EASD guide, Customer requests a cancellation to paid Order or OrderItems};
  \href{https://standards.atlassian.net/wiki/spaces/EASD/pages/842301447}{EASD guide, Customer requests a change of a paid order prior to travel resulting in additional collection and a refund};
  \href{https://standards.atlassian.net/wiki/spaces/EASD/pages/842268712}{EASD guide, Customer requests a change of a paid order prior to travel resulting in additional collection and a value store};
  \href{https://standards.atlassian.net/wiki/spaces/OA/pages/844169238}{Order Accounting, Customer requests a change of a paid order prior to travel resulting in additional collection and a value store}.

  Capabilities: ORDPEN, ORDCAN.
\end{theorem}

\begin{theorem}[Penalties and forfeits persist (C-95)]
  \label{thm:change-cancellation-penalties-persist}
  \uses{def:order, def:c18-change, def:money}
  From release 24.1, penalties and forfeited amounts are recorded at order
  level with links to order items. Once recorded, each keeps its amount and
  stays on the order when its order item is later replaced, cancelled or
  dropped from the order view.

  Case to check: In the change worked examples the change penalty points at
  the new item \texttt{ORDITM02\_NEW}, and the next retrieval no longer
  shows the cancelled item. A model that stores penalties on items would
  lose them at the next change.

  Sources: \href{https://retailing.iata.org/armi/docs/ORDPEN/}{IATA capability ORDPEN};
  \href{https://retailing.iata.org/armi/docs/ORDFOR/}{IATA capability ORDFOR};
  \href{https://standards.atlassian.net/wiki/spaces/EASD/pages/842268673}{EASD guide, Customer requests a change to an Order containing a paid OrderItem}.

  Capabilities: ORDPEN, ORDFOR.
\end{theorem}

\begin{definition}[Name change through an offer (C-96)]
  \label{def:change-cancellation-name-change}
  \uses{def:passenger, thm:delivered-not-cancellable, def:c08-delivery-status, def:c06-item-status}
  \notready
  A name change or correction is made by accepting a name change offer, free
  or paid. It adds a name change order item priced at the fee, whose service
  is \texttt{DELIVERED} at once; it overwrites all name elements of the same
  passenger and leaves the flight items and their services unchanged.
  Birthdate, gender, birthplace and additional names cannot be changed this
  way.

  Case to check: In the paid name change example the fee item ORDITM02 is
  \texttt{CLOSED} at 100 and the flight item ORDITM01 is unchanged. A full
  cancellation the next day would refund the flight item under its own rules
  and keep the 100.

  Open point: The source does not separate correcting a name from handing
  the booking to another person. It also does not say whether a paid name
  change fee is refunded if the order is later cancelled; because its
  service is \texttt{DELIVERED}, the existing rule forbids cancelling it.

  Sources: \href{https://retailing.iata.org/armi/docs/ORDNAM/}{IATA capability ORDNAM};
  \href{https://standards.atlassian.net/wiki/spaces/EASD/pages/672825346}{EASD guide, Customer Changes Their Name};
  \href{https://standards.atlassian.net/wiki/spaces/EASD/pages/672792581}{EASD guide, Paid name changes}.

  Capabilities: ORDNAM.
\end{definition}

\begin{definition}[Changes made without an offer (C-97)]
  \label{def:change-cancellation-no-offer-changes}
  \uses{def:c15-cancel-order, thm:c21-version, def:passenger, def:change-cancellation-partly-paid-cancel}
  \notready
  A change goes through an offer when it replaces, adds or cancels items of
  an order with any payment, removes a passenger, or changes a name. Other
  passenger details, sharing preferences, cancelling an unpaid order and
  accepting a follow-up action are made directly; they keep the passenger ID
  and have no payment step.

  Open point: The direct path has no price, yet the sample change record on
  the Modification of an Order page pairs an Accept follow-up action with
  100 due by the airline. The business must decide whether a direct change
  can create an amount owed and, if so, how it is priced.

  Sources: \href{https://standards.atlassian.net/wiki/spaces/EASD/pages/574554168}{EASD guide, Seller Initiated Change to an Order};
  \href{https://retailing.iata.org/armi/docs/ORDPAX/}{IATA capability ORDPAX};
  \href{https://standards.atlassian.net/wiki/spaces/EASD/pages/574619782}{EASD guide, Modification of an Order}.

  Capabilities: ORDPAX, ORDRSH, ORDNAM, ORDCAN, ORDCA2.
\end{definition}

\begin{theorem}[Full cancellation equals item cancellations (C-98)]
  \label{thm:change-cancellation-full-equals-items}
  \uses{def:c15-cancel-order, thm:c16-penalty-net, thm:change-cancellation-full-cancel-no-residual}
  \notready
  Cancelling a whole order produces the same refundable amounts and
  penalties as cancelling each of its items under the same rules.

  Case to check: An order has two items and a cancellation fee of 50 per
  transaction. Cancelling the order costs 50; cancelling the items one at a
  time costs 100.

  Open point: ORDCA2 says the calculation mechanisms are identical but does
  not say the totals must match. The business must decide whether
  order-level fees apply once per transaction or once per item, and whether
  item-by-item cancellation may keep reusable value that a full cancellation
  may not.

  Source: \href{https://retailing.iata.org/armi/docs/ORDCA2/}{IATA capability ORDCA2}.

  Capabilities: ORDCA2, ORDCAN.
\end{theorem}

\begin{theorem}[Stored value belongs to payer (C-99)]
  \label{thm:change-cancellation-stored-value-payer}
  \uses{def:payment, def:c13-several-methods, def:c17-refund-methods, thm:c20-value-store}
  \notready
  A value store is held by the airline for the payer of the transactions
  that created it.

  Open point: The source does not say whose value it is when several payers
  or payment methods paid the original item, or when the payer is not the
  passenger. The business must decide who may use, move or reclaim it.

  Source: \href{https://standards.atlassian.net/wiki/spaces/EASD/pages/574554245}{EASD guide, Value Store}.

  Capabilities: ORDREU.
\end{theorem}
